\section{Two arbitrary binary diameters}\label{sec:two}
We first justify attainment and separate the boundary cases from the
open stationary domain. Normalize $V=1$ and take a maximizing sequence
with ratio at least $c>0$. For two coordinates, $R\ge1/\sqrt2$,
whereas the discrepancy is at most one, so $R\le1/c$.
The squares are uniformly integrable, since
\[
 \sum_i\E[X_i^2;|X_i|>L]\le\frac1{cL}.
\]
After reflection, selected positive closed thresholds stay bounded:
Chebyshev controls thresholds tending to $-\infty$, and $1-\Phi(y)$
controls those tending to $+\infty$. Coordinatewise weak limits preserve
centering and total variance. A weak limit of two-point laws has at most
two points. Closed-set Portmanteau at the moving thresholds and lower
semicontinuity of $R$ give a limiting ratio at least the supremum;
equality follows from its definition. Thus a binary-class maximizer exists.

\subsection{The top threshold and equal-diameter boundary}
For a centered binary variable on $\{-b,a\}$, $a,b>0$,
\[
 v=ab,\qquad \rho=\frac{ab(a^2+b^2)}{a+b}
                 \ge2(\sqrt2-1)\frac{v^2}{a}.
\]
The inequality is the minimum of $(1+x^2)/(1+x)$ for $x=a/b>0$.
For two such coordinates with total upper endpoint $L$, Cauchy's
inequality gives $R\ge2(\sqrt2-1)V^2/L$.
At the top closed threshold the CDF is one, so its ratio is at most
\[
 \frac{z\{1-\Phi(z)\}}{2(\sqrt2-1)}
 \le\frac{e^{-1/2}}{4(\sqrt2-1)}<\frac{37}{100},
 \qquad z=L/\sqrt V.
\]
Here $1-\Phi(z)\le e^{-z^2/2}/2$ follows by translating its Gaussian
integral. For the final rational comparison,
$\sqrt e>41/25$, $\sqrt2-1>207/500$ give
$12500/33948<37/100$.
Reflection treats the negative strict bottom threshold.
The common-diameter two-active-bit case was excluded in
Section~\ref{sec:common}. A zero coordinate leaves exactly the
one-coordinate problem, whose ratio exceeds $37/100$.

\subsection{Open unequal-diameter stationary points}
Scale the larger diameter to one and write the other as $t\in(0,1)$.
Let their upper probabilities be $p,q\in(0,1)$, and put
\[
 a=p(1-p),\quad b=q(1-q),\quad
 r_p=p^2+(1-p)^2,\quad r_q=q^2+(1-q)^2,
 \quad V=a+t^2b,\quad R=ar_p+t^3br_q.
\]
The raw support is $0,t,1,1+t$. The three remaining positive closed
thresholds have the data
\[
 \begin{array}{c|c|c}
 \text{raw threshold}&F&y\ \text{(centered threshold)}\\\hline
 0&(1-p)(1-q)&-p-tq\\
 t&1-p&-p+t(1-q)\\
 1&1-pq&1-p-tq.
 \end{array}
\]
Set $z=y/\sqrt V$, $D=F-\Phi(z)$ and $C=DV^{3/2}/R$.
At a global maximum, a selected active branch has a local maximum
even if other thresholds tie. Its three derivatives therefore vanish.
Define
\[
 \alpha=\phi(z)/\sqrt V,\quad
 A=\frac{3D+z\phi(z)}{2V},\quad\lambda=D/R.
\]
After removing the common positive factor $V^{3/2}/R$, the two
probability derivatives are
\begin{align*}
 g_p&=F_p+\alpha+A(1-2p)-\lambda(1-2p)^3,\\
 g_q&=F_q+\alpha t+At^2(1-2q)-\lambda t^3(1-2q)^3.
\end{align*}
Divide the diameter derivative by $q$ for thresholds $0,1$ and by
$1-q$ for threshold $t$, then multiply by $VR\sqrt V$. Its numerator is
\begin{equation}\label{eq:gap-gradient}
 G=\phi(z)RB+3Dwt a\sqrt V\,(r_p-tr_q),
\end{equation}
where the pairs $(w,B)$ for these three rows are, respectively,
\[
 \bigl(1-q,\ p[(1-p)-t(1-q)]\bigr),\quad
 \bigl(q,-p[(1-p)+qt]\bigr),\quad
 \bigl(1-q,(1-p)[p+t(1-q)]\bigr).
\]
The extra factors $p$ and $1-p$ in the first and last rows are
strictly positive in the open chart and can also be removed.
These factorizations retain cancellation near vanishing probabilities.

\subsection{A compact domain for the necessary equations}
The probability equations also bound the normalized support points.
For one centered binary coordinate with law $P_i$, put
$s_i=\E X_i^2$, $b_i=\E|X_i|^3$, $M_i=\E X_i|X_i|$ and let
$G_i$ be the deleted CDF. In total-variance-one units define
\[
 Q_i(u)=G_i(y-u)-F(y)+\phi(y)u+A(u^2-s_i)
                 -C\{|u|^3-b_i-3M_i u\}.
\]
Its mean under $P_i$ is zero. Probability stationarity says its two
endpoint values agree, so both equal zero. This uses probability
stationarity alone, without a full-law majorant.
Since $Cb_i\le CR=D\le1$, $|M_i|\le s_i\le1$,
$\phi<2/5$ and $|A|\le13/8$, every active endpoint satisfies
\[
 C|u|^3\le4+(2/5+3C)|u|+(13/8)u^2.
\]
For $|u|\ge3$, dividing the right side by $u^2$ bounds it by
$C+793/360$. In particular $|u|<3+8/C$.
At $C\ge37/100$, the larger normalized diameter $1/\sqrt V$
is less than twice $911/37$. Hence
\begin{equation}\label{eq:two-compact-domain}
 V>\frac{1369}{3319684}>\frac1{2500},\qquad
 R\ge V^{3/2}/\sqrt2>V/75.
\end{equation}

The arithmetic certificate covers the three closed cubes
$(p,q,t)\in[0,1]^3$, conditional on \eqref{eq:two-compact-domain}
and $C\ge37/100$. Each box is excluded by an empty conditional domain,
a strict ratio upper bound, or a nonzero enclosure of a required
probability or factored diameter derivative. One useful ratio bound is
\[
 C\le D\sqrt{a/r_p^2+b/r_q^2},
\]
which follows from H\"older's inequality applied to $V,R$.
There are 37287 accepted source leaves in the complete partition.
An independently written, cancellation-preserving replay checks every
source condition, using 37421 descendant leaves and no unresolved
or unstarted region. The input bindings and exact binary forests are
part of the public certificate. Boundary boxes enclose limits of open
stationary points; they do not impose stationarity on actual one-bit
boundary maximizers.

Consequently no competitive unequal-diameter two-active-bit maximum
can satisfy its necessary equations. The preceding endpoint and
common-diameter arguments cover all remaining two-active-bit cases.
Attainment leaves exactly one active coordinate, proving
Theorem~\ref{thm:two}.
To require both variances positive, add a centered symmetric binary
coordinate of diameter tending to zero. Choose the closed threshold
just above the winning one-coordinate atom. Its CDF, variance and
third-moment denominator converge to the one-coordinate values.
Thus the supremum remains $C_1$, and the classification prevents its
attainment with both coordinates active.
